% META: {"id": "TEXP-CTP-EX-032", "chapitre": "TEXP-COMPLEXES-TRIGO-POLYNOMES", "type_objet": "exercice", "capacites": ["TEXP-CTP-C4"], "capacites_codes": ["C4"], "methodes": ["M4"], "parcours": 1, "competences": ["communiquer","raisonner"], "duree_min": 18, "mode_creation": "ex_nihilo", "sources_inspiration": [], "fichier_tex": "chapitres/TEXP-COMPLEXES-TRIGO-POLYNOMES/exercices/TEXP-CTP-EX-032.tex", "corrige_tex": "chapitres/TEXP-COMPLEXES-TRIGO-POLYNOMES/corriges/TEXP-CTP-CO-032.tex", "status": "approved"}
% BEGIN-VERIFY
% from sympy import *
% x = symbols('x')
% f = x*exp(-x)
% fp = diff(f, x)
% fpp = diff(f, x, 2)
% assert simplify(fp) == (1-x)*exp(-x)
% assert simplify(fpp) == (x-2)*exp(-x)
% assert solve(fp, x) == [1]
% assert solve(fpp, x) == [2]
% assert f.subs(x, 0) == 0
% assert f.subs(x, 1) == exp(-1)
% assert f.subs(x, 2) == 2*exp(-2)
% assert limit(f, x, oo) == 0
% END-VERIFY
\begin{exercice}{TEXP-CTP-EX-032}{2}{18}
Soit $f(x) = x\,\mathrm{e}^{-x}$ sur $\mathbb{R}$.
\begin{enumerate}
  \item Dresser le tableau de variations complet (limites, $f'$, $f''$).
  \item Identifier les points d'inflexion.
  \item Esquisser la courbe en placant l'asymptote et les points remarquables.
\end{enumerate}

\end{exercice}
